\clearpage
\section{Verdict response: ranks, abstention, and transfer boundaries}
\subsection{What the benchmark does and does not support}
The Ssym inverse comparison is the paper's positive benchmark, not a general license to predict biochemical rescue. The committed B1 bootstrap records 342 inverse mutations, inverse correlation 0.5919 and standard deviation of residuals 1.4292 kcal/mol. Its one-sided comparisons against published PoPMuSiCsym point values pass the locked gate, but method-paired prediction differences were not bootstrapped. The training/evaluation sequence background and historical comparator training overlap also require caution. It is a narrow benchmark outcome, not evidence that the predicted absolute energies on an unrelated disease protein are calibrated.

The external gradient contradicts a universal stability-prediction claim. Long training makes p53 transfer correlation fall from 0.426 to -0.02; validation-based early stopping does not restore it. The initial MYOGLOBIN correlation 0.455 included 41 same-protein training rows and falls to 0.24--0.29 on the corrected comparison. SOD1 measured stability gives $r=-0.15$ on eight mutations; MaveDB variant-abundance and activity comparisons are distinct phenotypes and cannot substitute for stability validation. The missed p53 Y220C suppressor A138G and the implausible +46 kcal/mol multi-mutation sum further delimit the scan. Absolute stability changes should not be inferred from rank positions.

\subsection{A decision tree for using or refusing a score}
\begin{enumerate}
\item Verify the exact deposited construct, background mutations, chain and residue numbering before scoring. The 2VUK chain includes a superstable quadruple-mutant background; treating it as wild-type p53 changes the counterfactual. If the mapping is ambiguous, abstain.
\item Ask whether the substitution and target belong to the model's observed support. The training set has no proline mutations, so proline-position scores are out-of-distribution and must be masked. For metal-dependent or oligomeric targets such as SOD1, a single-chain C$\alpha$ contact graph omits cofactor and interface mechanisms; do not use its rank as rescue evidence.
\item For a supported, single-chain soluble target with a verified background, use identity-bias-corrected ranks only as a triage list. Do not add single-mutant scores to assert a multi-mutant free-energy gain; the recorded +46 kcal/mol sum is a warning that additivity fails. Report uncertainty, source structure and the chosen checkpoint alongside each list.
\item Before making an activity, rescue or therapeutic claim, require an independent assay and a comparator on the same variants. Until then, even a top-ranked sequence is an unvalidated hypothesis. If the assay or comparator is missing, abstain from the biological claim.
\end{enumerate}

\subsection{What remains to test}
A contact-map or ESM-derived edge/feature ablation, metal-aware and dimer-aware SOD1 model, family-held-out cross-validation, an antisymmetry-consistent ProTherm benchmark, cross-species p53 transfer and a fresh preregistered rescue nomination are [PENDING]. These are proposed experiments, not repairs achieved by this manuscript revision. A family-level holdout needs proteins, homolog annotations and assay units frozen before training; row-level variants of the same protein are not independent families. Tool inventories must separate directly executed scientific tools from infrastructure and literature; the strict 40-tool gate remains open. This revision adds substantive manuscript pages and includes no slides.
